\documentclass[11pt]{article}

\usepackage[a4paper,margin=1in]{geometry}
\usepackage[T1]{fontenc}
\usepackage{lmodern}
\usepackage{amsmath,amssymb,amsthm,mathtools}
\usepackage{microtype}
\usepackage[hidelinks]{hyperref}
\hypersetup{
  pdftitle={A Dominant Polynomial Parametrization of xyz plus t squared plus 1 equals 0},
  pdfsubject={A polynomial parametrization of a Diophantine hypersurface},
  pdfkeywords={Diophantine equation, polynomial parametrization, Gaussian integers, norm identity}
}

\newtheorem{theorem}{Theorem}[section]
\newtheorem{lemma}[theorem]{Lemma}
\newtheorem{proposition}[theorem]{Proposition}
\newtheorem{corollary}[theorem]{Corollary}
\theoremstyle{remark}
\newtheorem{remark}[theorem]{Remark}

\newcommand{\A}{\mathbb A}
\newcommand{\Q}{\mathbb Q}
\newcommand{\Z}{\mathbb Z}

\title{A Dominant Polynomial Parametrization of\\
\texorpdfstring{$xyz+t^2+1=0$}{xyz+t^2+1=0}}
\author{}
\date{}

\begin{document}
\maketitle
\vspace{-2.2em}

\begin{abstract}
We exhibit explicit polynomials $X,Y,Z,T\in\Z[a,b,c]$ satisfying
$XYZ+T^2+1=0$ identically.  The construction comes from an iterated
Gaussian-norm factorization that preserves imaginary part $1$.  The
resulting three-parameter polynomial map is dominant on the
hypersurface, and a nonconstant one-parameter specialization is also
recorded.
\end{abstract}

\section{Main result}

\begin{theorem}\label{thm:main}
Let $a,b,c$ be independent parameters, put
\begin{equation}\label{eq:R-def}
R=a+b(a^2+1),
\end{equation}
and define
\begin{equation}\label{eq:parametrization}
\boxed{
\begin{aligned}
Z&=a^2+1,\\
Y&=(1+ab)^2+b^2,\\
T&=R+c(R^2+1),\\
X&=-\bigl((1+cR)^2+c^2\bigr).
\end{aligned}}
\end{equation}
Then $X,Y,Z,T\in\Z[a,b,c]$ and
\begin{equation}\label{eq:main-identity}
XYZ+T^2+1=0
\end{equation}
identically.  Hence every $(a,b,c)\in\Z^3$ gives an integer solution
$(x,y,z,t)=(X,Y,Z,T)$ of $xyz+t^2+1=0$.  For every such specialization,
\[
X<0,\qquad Y>0,\qquad Z>0.
\]
\end{theorem}

Eliminating the auxiliary symbol $R$, the map is explicitly
\begin{equation}\label{eq:expanded-parametrization}
\begin{aligned}
X={}&-\left(\left(1+c\bigl(a+b(a^2+1)\bigr)\right)^2+c^2\right),\\
Y={}&(1+ab)^2+b^2,\\
Z={}&a^2+1,\\
T={}&a+b(a^2+1)
+c\left(\bigl(a+b(a^2+1)\bigr)^2+1\right).
\end{aligned}
\end{equation}

\section{Proof}

The following identity is the algebraic mechanism behind the
parametrization.

\begin{lemma}[Norm splitting]\label{lem:norm-splitting}
For arbitrary elements $U,V$ of a commutative ring, put
\[
S=U+V(U^2+1).
\]
Then
\begin{equation}\label{eq:norm-identity}
S^2+1=(U^2+1)\bigl((1+UV)^2+V^2\bigr).
\end{equation}
\end{lemma}

\begin{proof}
In the universal polynomial ring with an element $i$ satisfying
$i^2=-1$, one has
\begin{equation}\label{eq:gaussian-factorization}
S+i=(U+i)(1+UV-Vi).
\end{equation}
The real part of the product on the right is
\[
U(1+UV)+V=U+V(U^2+1)=S,
\]
and its imaginary part is $(1+UV)-UV=1$.  Multiplication by the
conjugate of \eqref{eq:gaussian-factorization} gives
\eqref{eq:norm-identity}.  Since this is a polynomial identity over
$\Z$, it remains valid after specialization to any commutative ring.
\end{proof}

\begin{proof}[Proof of Theorem~\ref{thm:main}]
Apply Lemma~\ref{lem:norm-splitting} first with $(U,V)=(a,b)$.  By
\eqref{eq:R-def} and \eqref{eq:parametrization},
\begin{equation}\label{eq:first-splitting}
R^2+1=(a^2+1)\bigl((1+ab)^2+b^2\bigr)=ZY.
\end{equation}
Apply the lemma a second time with $(U,V)=(R,c)$.  Since
$T=R+c(R^2+1)$,
\begin{align*}
T^2+1
  &=(R^2+1)\bigl((1+cR)^2+c^2\bigr)\\
  &=ZY\bigl((1+cR)^2+c^2\bigr)\\
  &=-XYZ.
\end{align*}
This proves \eqref{eq:main-identity}.

For integral $a,b,c$, one has $Z=a^2+1\geq1$.  The expression
$Y=(1+ab)^2+b^2$ cannot vanish, since $b=0$ would leave the first
square equal to $1$; hence $Y\geq1$.  The same argument shows that
$(1+cR)^2+c^2\geq1$, so $X\leq-1$.
\end{proof}

\section{How the construction was found}\label{sec:discovery}

Writing the equation as
\begin{equation}\label{eq:target-factorization}
t^2+1=-xyz
\end{equation}
shows that one needs to split the norm $t^2+1=N(t+i)$ into three
polynomial factors.  Multiplication by a general Gaussian integer does
not preserve the special form ``real part plus $i$''.  The key is to
choose a multiplier whose imaginary cross-terms cancel.

Starting from $U+i$, multiply by $1+UV-Vi$.  Then
\[
(U+i)(1+UV-Vi)=\bigl(U+V(U^2+1)\bigr)+i.
\]
Thus the imaginary part remains exactly $1$, while taking norms gives
\[
\bigl(U+V(U^2+1)\bigr)^2+1
=(U^2+1)\bigl((1+UV)^2+V^2\bigr).
\]
The output is again a norm of the same shape, so the operation can be
iterated.

The first iteration, with $(U,V)=(a,b)$, produces
\[
R^2+1=ZY.
\]
The second, with $(U,V)=(R,c)$, produces
\[
T^2+1=(R^2+1)\bigl((1+cR)^2+c^2\bigr)
      =ZY\bigl((1+cR)^2+c^2\bigr).
\]
Taking $X$ to be the negative of the final factor gives
$T^2+1=-XYZ$.  The parametrization is therefore a two-stage
composition of Gaussian-norm splittings, rather than an isolated
polynomial expansion.

\section{Dominance and scope}

Let
\[
\mathcal H=\{(x,y,z,t)\in\A^4_{\Q}:xyz+t^2+1=0\}.
\]
Theorem~\ref{thm:main} defines a polynomial map
\[
\Phi:\A^3_{\Q}\longrightarrow\mathcal H,
\qquad
(a,b,c)\longmapsto(X,Y,Z,T).
\]

\begin{proposition}\label{prop:dominant}
The map $\Phi$ is dominant; equivalently, its image is Zariski dense
in $\mathcal H$.
\end{proposition}

\begin{proof}
Consider the coordinate functions $(Z,Y,T)$.  From
\eqref{eq:parametrization} and \eqref{eq:first-splitting},
\[
\frac{\partial Z}{\partial a}=2a,
\qquad
\frac{\partial Y}{\partial b}=2R,
\qquad
\frac{\partial T}{\partial c}=R^2+1=YZ.
\]
The relevant Jacobian is therefore
\begin{equation}\label{eq:jacobian}
J_{\Phi}:=
\det\left(\frac{\partial(Z,Y,T)}{\partial(a,b,c)}\right)
=
\det
\begin{pmatrix}
2a&0&0\\
\partial_aY&2R&0\\
\partial_aT&\partial_bT&YZ
\end{pmatrix}
=4aRYZ.
\end{equation}
It is not the zero polynomial; at $(a,b,c)=(1,0,0)$ it equals $8$.
Hence $\Phi$ has generic rank $3$, and the closure of its image has
dimension $3$.

The polynomial $xyz+t^2+1$ is irreducible over $\Q$.  Indeed, as a
polynomial in $x$ over the unique factorization domain $\Q[y,z,t]$, it
is primitive because $yz$ and $t^2+1$ are coprime.  Over the fraction
field $\Q(y,z,t)$ it is linear and therefore irreducible; Gauss's
lemma gives irreducibility over $\Q[y,z,t]$.  Thus $\mathcal H$ is an
irreducible hypersurface of dimension $3$.  The Zariski closure of
$\Phi(\A^3_{\Q})$ is irreducible and has dimension $3$; because it is
contained in $\mathcal H$, it must equal $\mathcal H$.  This proves
dominance.
\end{proof}

\begin{remark}\label{rem:scope}
Dominance says that the construction is full-dimensional, not merely
a curve or surface in $\mathcal H$.  It does not claim that every
integer point of $\mathcal H$ is attained by integral parameters.
\end{remark}

\section{A one-parameter specialization}

Taking $a=n$, $b=-1$, and $c=1$ gives a family in which all four
coordinates are nonconstant.

\begin{corollary}\label{cor:one-parameter}
For every $n\in\Z$, set
\begin{equation}\label{eq:one-param}
\boxed{
\begin{aligned}
x&=-\bigl((n^2-n)^2+1\bigr),\\
y&=(n-1)^2+1,\\
z&=n^2+1,\\
t&=n^4-2n^3+2n^2-n+1.
\end{aligned}}
\end{equation}
Then $xyz+t^2+1=0$.
\end{corollary}

\begin{proof}
Here
\[
R=n-(n^2+1)=-n^2+n-1.
\]
Substitution into \eqref{eq:parametrization} gives
\eqref{eq:one-param}, and Theorem~\ref{thm:main} applies.  Equivalently,
\begin{align*}
&\bigl(n^4-2n^3+2n^2-n+1\bigr)^2+1\\
&\qquad={}
\bigl((n^2-n)^2+1\bigr)
\bigl((n-1)^2+1\bigr)
\bigl(n^2+1\bigr).
\end{align*}
\end{proof}

For example, $n=2$ gives $(x,y,z,t)=(-5,2,5,7)$, and
\[
(-5)(2)(5)+7^2+1=0.
\]

\end{document}
